%============================================================
% CH1 · P2 · PROBLEM 2 — REMOVE DUPLICATES FROM SORTED ARRAY II (LeetCode 80)
%============================================================
\subsection{Problem 1.2.2: Remove Duplicates from Sorted Array II}
\problemlabel{1.2.2}{ssec:ch1-remove-duplicates-ii}

\problemcard{1.2.2}{Remove Duplicates from Sorted Array II}{LeetCode 80}{https://leetcode.com/problems/remove-duplicates-from-sorted-array-ii/}{\exmedium\ \lc{Medium}}{Read/write pointers with a rule on the written part: keep $x$ iff $x \ne a_{w-2}$}

\subsubsection*{1. Problem statement}

An integer array \texttt{nums} is sorted in non-decreasing order. Remove duplicates \textbf{in place} so that every distinct value appears \textbf{at most twice}, keeping the relative order. Return the number $k$ of elements that remain; the first $k$ cells of \texttt{nums} must hold them, and what is stored beyond the first $k$ cells does not matter. Only $\bigO{1}$ extra memory is allowed.

\textbf{Constraints:} $1 \le n \le 3\cdot10^4$, \; $-10^4 \le \texttt{nums}[i] \le 10^4$, \; \texttt{nums} is sorted in non-decreasing order.

\begin{table}[H]
\centering\small
\begin{tabular}{@{}llcl@{}}
\toprule
& \textbf{nums} & $k$ & \textbf{nums[0..k)} \\
\midrule
Example 1 & $[1, 1, 1, 2, 2, 3]$ & 5 & $[1, 1, 2, 2, 3]$ \\
Example 2 & $[0, 0, 1, 1, 1, 1, 2, 3, 3]$ & 7 & $[0, 0, 1, 1, 2, 3, 3]$ \\
\bottomrule
\end{tabular}
\end{table}

\begin{diagrambox}
\begin{tikzpicture}
    \begin{scope}
        \dparray{e}{0}{Ex.\ 1}{1/akeep, 1/akeep, 1/aout, 2/akeep, 2/akeep, 3/akeep}
        \node[note, anchor=west] at (5.6,0) {$k = 5$};
    \end{scope}
    \begin{scope}[yshift=-1.0cm]
        \dparray{f}{0}{Ex.\ 2}{0/akeep, 0/akeep, 1/akeep, 1/akeep, 1/aout, 1/aout, 2/akeep, 3/akeep, 3/akeep}
        \node[note, anchor=west] at (8.3,0) {$k = 7$};
    \end{scope}
\end{tikzpicture}
\figcap{Green elements survive, grey ones are removed: from every run of equal values only the first two copies stay.}
\end{diagrambox}

\subsubsection*{2. The problem model}
\begin{modelbox}[Problem Model --- Remove Duplicates from Sorted Array II]
\renewcommand{\arraystretch}{1.35}
\begin{tabularx}{\linewidth}{@{}>{\bfseries}l X@{}}
Input & a \emph{sorted} array: equal values form contiguous runs \\
Computed & a length $k$ and, in place, the subsequence that keeps $\min(2, \text{run length})$ copies of every run \\
Over & elements; stable compaction (Pattern 2) with a rule that depends on the already kept elements \\
Sizes & $n \le 3\cdot10^4$: $\bigO{n^2} \approx 9\cdot10^8$ shifts in the worst case is too slow; aim for $\bigO{n}$ time, $\bigO{1}$ memory \\
Key fact & sortedness turns ``already two copies of $x$ kept?'' into one comparison \\
\end{tabularx}
\end{modelbox}

\begin{diagrambox}
\begin{tikzpicture}
    \dparray{m}{0}{nums}{0/awin, 0/awin, 1/akeep, 1/akeep, 1/akeep, 1/akeep, 2/awin, 3/akeep, 3/akeep}
    \arange{m0}{m1}{keep 2}{memoorange}
    \arange{m2}{m5}{keep 2 of 4}{successgreen}
    \arange{m6}{m6}{keep 1}{memoorange}
    \arange{m7}{m8}{keep 2}{successgreen}
\end{tikzpicture}
\figcap{The model in terms of runs: every run of length $\ell$ contributes $\min(2, \ell)$ elements, so $k = 2 + 2 + 1 + 2 = 7$ for Example 2.}
\end{diagrambox}

\subsubsection*{3. Key observations}

\begin{itemize}[leftmargin=*]
    \item \textbf{The kept part is sorted too.} Compaction preserves order, so \texttt{nums[0..w)} is sorted. If two copies of $x$ have already been kept, they are the last two kept elements, $a_{w-2} = a_{w-1} = x$.
    \item \textbf{One comparison is enough.} Keep $x$ iff $w < 2$ or $x \ne a_{w-2}$. Indeed, if $a_{w-2} = x$, then $a_{w-2} \le a_{w-1} \le x$ forces $a_{w-1} = x$ as well: two copies are already there. If $a_{w-2} \ne x$ (it must be smaller), at most one copy of $x$ was kept.
    \item \textbf{Compare with the written part, not the read part.} The tempting rule $x \ne a_{r-2}$ looks at a cell that may have been overwritten. On Example 1 it returns $k = 4$ with $[1, 1, 2, 3]$: when the second 2 is read ($r = 4$), cell $a_2$ already holds the first 2, so the second 2 is mistaken for a third copy.
    \item \textbf{$k = 1$ is the classic version.} With ``at most once'' the rule is $x \ne a_{w-1}$ (LeetCode 26, variant V1); with ``at most $m$'' it is $x \ne a_{w-m}$.
\end{itemize}

\begin{diagrambox}
\begin{tikzpicture}
    \begin{scope}
        \dparray{x}{0}{at $r = 4$}{1/akeep, 1/akeep, 2/awrite, 2/acell, 2/aread, 3/acell}
        \aptrup{x2}{a[r-2]}{warningred}
        \aptr{x4}{r}{successgreen}
        \aptr{x3}{w}{accentorange}
        \node[note, anchor=west, text width=4.2cm, align=left] at (5.8,0) {cell 2 was overwritten at $r = 3$ (it held a 1); now it shows a 2 and the rule $x \ne a_{r-2}$ wrongly drops the second 2};
    \end{scope}
\end{tikzpicture}
\figcap{Why the rule must use $w$: a cell between $w$ and $r$ no longer holds its original value. The correct rule compares with $a_{w-2} = a_1 = 1$ and keeps the 2.}
\end{diagrambox}

\subsubsection*{4. Approaches}

\begin{bruteforcebox}
\textbf{Idea.} Walk with an index $i$ over the current (shrinking) array; whenever \texttt{nums[i]} equals \texttt{nums[i-2]} it is a third copy: delete it by shifting the whole tail one step left. In place, but every deletion costs $\bigO{n}$.
\begin{lstlisting}[style=cpp]
#include <bits/stdc++.h>
using namespace std;

class Solution {
public:
    int removeDuplicates(vector<int>& nums) {
        int len = nums.size(), i = 2;
        while (i < len) {
            if (nums[i] == nums[i - 2]) {                    // a third copy: delete it
                for (int j = i; j + 1 < len; j++) nums[j] = nums[j + 1];
                len--;                                       // do not advance i
            } else i++;
        }
        return len;
    }
};

int main() {
    vector<vector<int>> tests = {{1, 1, 1, 2, 2, 3}, {0, 0, 1, 1, 1, 1, 2, 3, 3}};
    for (auto& t : tests) {
        int k = Solution().removeDuplicates(t);
        cout << k << ":";
        for (int i = 0; i < k; i++) cout << " " << t[i];   // 5: 1 1 2 2 3 / 7: 0 0 1 1 2 3 3
        cout << "\n";
    }
}
\end{lstlisting}
\begin{complexitybox}
\textbf{Time:} $\bigO{n^2}$ in the worst case --- an array of $3\cdot10^4$ equal values deletes $n - 2$ elements, each shifting up to $n$ cells: about $4.5\cdot10^8$ moves. \textbf{Space:} $\bigO{1}$. Here comparing with \texttt{nums[i-2]} \emph{is} correct, because the array really is shifted: positions $< i$ are final.
\end{complexitybox}
\end{bruteforcebox}

\begin{algorithmbox}
\textbf{Idea.} Pattern 2 with the keep rule of the observations: one read pass; an element $x$ is written to \texttt{nums[w++]} iff $w < 2$ or $x \ne \texttt{nums[w-2]}$. The answer is $w$.
\begin{lstlisting}[style=cpp]
#include <bits/stdc++.h>
using namespace std;

class Solution {
public:
    int removeDuplicates(vector<int>& nums) {
        int w = 0;                                      // nums[0..w) = kept elements
        for (int x : nums)                              // x = nums[r], read before any write to r
            if (w < 2 || x != nums[w - 2]) nums[w++] = x;   // fewer than 2 copies kept so far
        return w;
    }
};

int main() {
    vector<vector<int>> tests = {{1, 1, 1, 2, 2, 3}, {0, 0, 1, 1, 1, 1, 2, 3, 3}, {1}, {1, 1, 1}};
    for (auto& t : tests) {
        int k = Solution().removeDuplicates(t);
        cout << k << ":";
        for (int i = 0; i < k; i++) cout << " " << t[i];
        cout << "\n";                       // 5: 1 1 2 2 3 / 7: 0 0 1 1 2 3 3 / 1: 1 / 2: 1 1
    }
}
\end{lstlisting}
\begin{complexitybox}
\textbf{Time:} $\bigO{n}$ --- one comparison and at most one write per element. \textbf{Space:} $\bigO{1}$. \textbf{Verdict:} \ac{}.
\end{complexitybox}
\end{algorithmbox}

The range-\texttt{for} reads \texttt{x} from \texttt{nums[r]} before the body may write \texttt{nums[w]} with $w \le r$; writing to the vector does not invalidate the loop because its size never changes.

\begin{invariantbox}[Why it is correct]
\textbf{Invariant.} Before reading index $r$: $w \le r$, and \texttt{nums[0..w)} is the required output for the prefix \texttt{nums[0..r)}, i.e.\ each run of that prefix contributes its first $\min(2, \ell)$ copies, in order. Moreover \texttt{nums[0..w)} is sorted.

\textbf{Maintenance.} Let $x = \texttt{nums}[r]$ (still original, since $w \le r$). In the sorted prefix, $x$ is at least every earlier value, so the copies of $x$ already kept are the last elements of \texttt{nums[0..w)}.
\begin{itemize}[leftmargin=*]
    \item If $w \ge 2$ and $\texttt{nums}[w-2] = x$: then $\texttt{nums}[w-2] \le \texttt{nums}[w-1] \le x$ gives $\texttt{nums}[w-1] = x$ too. Two copies are kept; $x$ is a third or later copy and must be dropped --- the code drops it.
    \item Otherwise ($w < 2$, or $\texttt{nums}[w-2] < x$): at most one copy of $x$ is among the kept elements, so $x$ must be kept --- the code appends it at position $w \le r$, a cell already read. The kept part stays sorted because $x \ge \texttt{nums}[w-1]$.
\end{itemize}
\textbf{Termination.} At $r = n$, \texttt{nums[0..w)} is the required output for the whole array, and the function returns its length. \hfill$\blacksquare$
\end{invariantbox}

\begin{alternativebox}
\textbf{Run by run.} Use the run structure directly: from the read position $r$, find the end $e$ of the run of \texttt{nums[r]} (the first index with a different value), write $\min(2, e - r)$ copies of the value, and jump to $r = e$. This is the ``group pointer'' that Chapter~\ref{ch:two_pointers} develops; it generalises to rules that depend on the whole run (``keep runs of length at least 3'', ``replace each run by its length'').
\begin{lstlisting}[style=cpp]
#include <bits/stdc++.h>
using namespace std;

class Solution {
public:
    int removeDuplicates(vector<int>& nums) {
        int n = nums.size(), w = 0, r = 0;
        while (r < n) {
            int e = r;
            while (e < n && nums[e] == nums[r]) e++;     // the run is [r, e)
            int keep = min(2, e - r);
            int v = nums[r];                             // save: the writes may reach r
            for (int c = 0; c < keep; c++) nums[w++] = v;
            r = e;                                       // jump to the next run
        }
        return w;
    }
};

int main() {
    vector<vector<int>> tests = {{1, 1, 1, 2, 2, 3}, {0, 0, 1, 1, 1, 1, 2, 3, 3}};
    for (auto& t : tests) {
        int k = Solution().removeDuplicates(t);
        cout << k << ":";
        for (int i = 0; i < k; i++) cout << " " << t[i];   // 5: 1 1 2 2 3 / 7: 0 0 1 1 2 3 3
        cout << "\n";
    }
}
\end{lstlisting}
Every index is visited once by $e$, so the time is $\bigO{n}$; memory $\bigO{1}$. The writes never pass the read position: before handling a run, $w \le r$, and the run writes at most $\min(2, e - r) \le e - r$ cells, so afterwards $w \le e$.

\begin{diagrambox}
\begin{tikzpicture}
    \dparray{g}{0}{nums}{0/akeep, 0/akeep, 1/akeep, 1/akeep, 1/aout, 1/aout, 2/aread, 3/acell, 3/acell}
    \aptr{g4}{w}{accentorange}
    \aptr{g6}{r}{successgreen}
    \aptrup{g7}{e}{primaryblue}
    \node[note, anchor=west, text width=3.6cm, align=left] at (8.3,0) {run $[6, 7)$ of value 2 has length 1: write one copy at $w = 4$};
\end{tikzpicture}
\figcap{Example 2 after the runs of 0's and 1's: four elements kept, the run of the 2 is next.}
\end{diagrambox}
\end{alternativebox}

\begin{tricksbox}
\begin{itemize}[leftmargin=*]
    \item \textbf{Any limit $m$:} replace both 2's by $m$ --- \texttt{if (w < m || x != nums[w - m])}. For $m = 1$ this is the classic deduplication (variant V1).
    \item \textbf{Why the array must be sorted:} on an unsorted array such as $[1, 2, 1, 1]$ the third 1 has $a_{w-2} = 2 \ne 1$ and would be kept. Unsorted inputs need counts (Pattern~5) --- variant V4.
    \item \textbf{A counter version} (\texttt{cnt} of the current value in the kept part) is also correct but has more branches; the $a_{w-2}$ rule hides the counter inside the array itself.
\end{itemize}
\end{tricksbox}

\subsubsection*{5. Trace}

\begin{tracebox}[{Example 2: $\texttt{nums} = [0, 0, 1, 1, 1, 1, 2, 3, 3]$}]
\centering
\begin{tikzpicture}
    \begin{scope}
        \dparray{a}{0}{$r = 3$}{0/akeep, 0/akeep, 1/akeep, 1/awrite, 1/acell, 1/acell, 2/acell, 3/acell, 3/acell}
        \aptr{a3}{r}{successgreen}
        \node[note, anchor=west, text width=2.9cm, align=left] at (8.3,0) {$a_{w-2} = a_1 = 0 \ne 1$: keep; $w = 4$};
    \end{scope}
    \begin{scope}[yshift=-1.75cm]
        \dparray{b}{0}{$r = 5$}{0/akeep, 0/akeep, 1/akeep, 1/akeep, 1/aout, 1/aread, 2/acell, 3/acell, 3/acell}
        \aptr{b5}{r}{successgreen}
        \aptr{b4}{w}{accentorange}
        \node[note, anchor=west, text width=2.9cm, align=left] at (8.3,0) {$a_2 = 1 = x$: drop (fourth copy); $w$ stays 4};
    \end{scope}
    \begin{scope}[yshift=-3.5cm]
        \dparray{c}{0}{$r = 6$}{0/akeep, 0/akeep, 1/akeep, 1/akeep, 2/awrite, 1/aout, 2/aread, 3/acell, 3/acell}
        \aptr{c6}{r}{successgreen}
        \node[note, anchor=west, text width=2.9cm, align=left] at (8.3,0) {$a_2 = 1 \ne 2$: keep, written to $a_4$; $w = 5$};
    \end{scope}
    \begin{scope}[yshift=-5.25cm]
        \dparray{d}{0}{end}{0/adone, 0/adone, 1/adone, 1/adone, 2/adone, 3/adone, 3/adone, 3/aout, 3/aout}
        \node[note, anchor=west, text width=2.9cm, align=left] at (8.3,0) {$k = 7$; cells 7 and 8 hold leftovers};
    \end{scope}
\end{tikzpicture}
\figcap{Selected steps. Blue: the cell written in this step; grey: free cells between $w$ and $r$.}

\medskip
\small
\begin{tabular}{@{}cccccl@{}}
\toprule
$r$ & $x$ & $w$ before & \textbf{test} & \textbf{action} & \textbf{nums after the step} \\
\midrule
0 & 0 & 0 & $w < 2$ & keep & $[0, 0, 1, 1, 1, 1, 2, 3, 3]$ \\
1 & 0 & 1 & $w < 2$ & keep & $[0, 0, 1, 1, 1, 1, 2, 3, 3]$ \\
2 & 1 & 2 & $a_0 = 0 \ne 1$ & keep & $[0, 0, 1, 1, 1, 1, 2, 3, 3]$ \\
3 & 1 & 3 & $a_1 = 0 \ne 1$ & keep & $[0, 0, 1, 1, 1, 1, 2, 3, 3]$ \\
4 & 1 & 4 & $a_2 = 1 = x$ & drop & $[0, 0, 1, 1, 1, 1, 2, 3, 3]$ \\
5 & 1 & 4 & $a_2 = 1 = x$ & drop & $[0, 0, 1, 1, 1, 1, 2, 3, 3]$ \\
6 & 2 & 4 & $a_2 = 1 \ne 2$ & keep & $[0, 0, 1, 1, \mathbf{2}, 1, 2, 3, 3]$ \\
7 & 3 & 5 & $a_3 = 1 \ne 3$ & keep & $[0, 0, 1, 1, 2, \mathbf{3}, 2, 3, 3]$ \\
8 & 3 & 6 & $a_4 = 2 \ne 3$ & keep & $[0, 0, 1, 1, 2, 3, \mathbf{3}, 3, 3]$ \\
\bottomrule
\end{tabular}
\end{tracebox}

\begin{tracebox}[{Example 1, correct rule versus the $a_{r-2}$ rule}]
\centering\footnotesize\setlength{\tabcolsep}{4pt}
\begin{tabular}{@{}cc|ccl|ccl@{}}
\toprule
& & \multicolumn{3}{c|}{\textbf{rule $x \ne a_{w-2}$}} & \multicolumn{3}{c}{\textbf{rule $x \ne a_{r-2}$ (wrong)}} \\
$r$ & $x$ & \textbf{test} & \textbf{act} & \textbf{array} & \textbf{test} & \textbf{act} & \textbf{array} \\
\midrule
2 & 1 & $a_0 = 1$ & drop & $[1,1,1,2,2,3]$ & $a_0 = 1$ & drop & $[1,1,1,2,2,3]$ \\
3 & 2 & $a_0 = 1$ & keep & $[1,1,\mathbf{2},2,2,3]$ & $a_1 = 1$ & keep & $[1,1,\mathbf{2},2,2,3]$ \\
4 & 2 & $a_1 = 1$ & keep & $[1,1,2,\mathbf{2},2,3]$ & $a_2 = \mathbf{2}$ & drop & $[1,1,2,2,2,3]$ \\
5 & 3 & $a_2 = 2$ & keep & $[1,1,2,2,\mathbf{3},3]$ & $a_3 = 2$ & keep & $[1,1,2,\mathbf{3},2,3]$ \\
\midrule
\multicolumn{2}{c|}{result} & \multicolumn{3}{l|}{$k = 5$: $[1, 1, 2, 2, 3]$} & \multicolumn{3}{l}{$k = 4$: $[1, 1, 2, 3]$ \wa{}} \\
\bottomrule
\end{tabular}
\end{tracebox}

\subsubsection*{6. Test suite}

\begin{testbox}
\centering\small
\begin{tabular}{@{}llc@{}}
\toprule
\textbf{nums} & \textbf{Why it is interesting} & \textbf{Expected $k$, prefix} \\
\midrule
$[1]$ & a single element & 1, $[1]$ \\
$[1, 1]$ & exactly two copies & 2, $[1, 1]$ \\
$[1, 1, 1]$ & smallest input with a removal & 2, $[1, 1]$ \\
$[1, 2, 3]$ & no duplicates at all & 3, $[1, 2, 3]$ \\
$[1, 1, 2, 2, 2, 2, 2, 3]$ & a long run in the middle & 5, $[1, 1, 2, 2, 3]$ \\
$[-10^4 \times 3, 10^4 \times 3]$ & extreme values, two runs of three & 4, $[-10^4, -10^4, 10^4, 10^4]$ \\
$[1, 1, 1, 2, 2, 3]$ & official Example 1 ($a_{r-2}$ trap) & 5, $[1, 1, 2, 2, 3]$ \\
$[0, 0, 1, 1, 1, 1, 2, 3, 3]$ & official Example 2 & 7, $[0, 0, 1, 1, 2, 3, 3]$ \\
$[7 \times 30000]$ & largest $n$, all equal (slow brute force) & 2, $[7, 7]$ \\
$[0, 0, 1, 1, \dots, 14999, 14999]$ & largest $n$, nothing to remove & 30000 \\
\bottomrule
\end{tabular}
\end{testbox}
The main solution, the run-by-run version and the in-place brute force reproduce both official examples and agree with a map-based reference (count copies, keep the first two) on 900 random sorted arrays ($n \le 9$, values in $[-3, 3]$).

\subsubsection*{7. Common mistakes}

\begin{warningbox}[Mistake 1 --- comparing with the read side]
\begin{lstlisting}[style=snippet]
if (r < 2 || nums[r] != nums[r - 2]) nums[w++] = nums[r];   // WRONG: Example 1 gives k = 4
if (w < 2 || nums[r] != nums[w - 2]) nums[w++] = nums[r];   // RIGHT
\end{lstlisting}
\end{warningbox}

\begin{warningbox}[Mistake 2 --- comparing with $a_{w-1}$ only]
\texttt{x != nums[w-1]} keeps one copy per value: that is LeetCode 26, not this problem. To allow two copies you must look two cells back.
\end{warningbox}

\begin{warningbox}[Mistake 3 --- \texttt{erase} on the vector]
Erasing third copies with \texttt{nums.erase} is the brute force in disguise: $\bigO{n^2}$ on an array of equal values, and easy to get wrong with the loop index.
\end{warningbox}

\begin{warningbox}[Mistake 4 --- returning the array instead of $k$]
The judge reads only \texttt{nums[0..k)}. Returning \texttt{nums.size()}, or forgetting that the tail still contains old values, prints garbage after the correct prefix.
\end{warningbox}

\subsubsection*{8. Takeaways}

\begin{itemize}[leftmargin=*]
    \item The keep rule of a read/write loop may look at the \emph{written} part \texttt{nums[0..w)} --- it is final --- but never at the cells between $w$ and $r$.
    \item Sorted input $+$ ``at most $m$ copies'' $=$ one comparison $x \ne a_{w-m}$.
    \item Thinking in runs (group pointer) is an equivalent, more general view: it handles any rule that depends on the length of a run.
    \item The worst case for the brute force is the ``boring'' input (all equal): always test it.
\end{itemize}

\begin{variantsbox}
\variant{At most once}{\exeasy\ \lc{Easy}\ \varlink{Remove Duplicates from Sorted Array}{LeetCode 26}{https://leetcode.com/problems/remove-duplicates-from-sorted-array/}}
  {keep every distinct value exactly once.}
  {$m = 1$: the kept part is the set of values in increasing order.}
  {keep $x$ iff $w < 1$ or $x \ne a_{w-1}$; $[1,1,2] \to 2$, $[0,0,1,1,1,2,2,3,3,4] \to 5$.}
\variant{At most $m$ copies}{\varnone}
  {every value may appear at most $m$ times ($m$ given).}
  {the same model with $\min(m, \ell)$ copies per run.}
  {keep $x$ iff $w < m$ or $x \ne a_{w-m}$; $m = 2$ is this problem.}
\variant{Only values that are not duplicated}{\varnone}
  {remove \emph{every} copy of a value that appears more than once; keep the values that occur exactly once.}
  {a decision per run that depends on its length, not on the last kept cells.}
  {run-by-run (Approach 3): write the value iff $e - r = 1$; $[1,1,2,3,3,4] \to [2,4]$.}
\variant{Unsorted array}{\varnone}
  {as this problem, but the array is not sorted and the order must be kept.}
  {copies of a value are no longer adjacent: the $a_{w-2}$ test is wrong ($[1,2,1,1]$).}
  {a frequency array of kept copies (Pattern~5) or a hash map, then the same compaction with \texttt{cnt[x] < 2}.}
\end{variantsbox}
